\documentclass[book.tex]{subfiles}

\begin{document}

\chapter{Equivariant Neural Networks}
\label{chap:equivariant_networks}
\noindent\rule{11cm}{0.4pt} \\
\textbf{Prerequisites:}  \autoref{chap:ml_basics}, \autoref{chap:grad_descent}  \\
\textbf{Difficulty Level:} ** \\
\noindent\rule{11cm}{0.4pt}
\vspace{5mm}

In this chapter we define what it means for a neural network to be equivariant to a group $G$. This definition will prove useful as we define concrete equivariant networks in subsequent chapters.

\section{$G$-function Transformations}
Suppose that $\mathcal{X}$ is a space upon which group $G$ acts. An element $g \in G$ has an action on functions
\begin{align}
f: \mathcal{X} \to \mathbb{R}^n
\end{align}
by the definition
\begin{align}
g \cdot f (x) &= f(g^{-1}x)
\end{align}
The transformation
\begin{align}
f \mapsto g \cdot f
\end{align}
is sometimes referred to as a $G$-function transform.


\section{Equivariant Neural Layers}


We can define a neural network as a function
\begin{align}
\psi_\theta: \mathcal{V} \to \mathcal{Y}
\end{align}
where $\theta$ is a set of weight parameters. Note here that \begin{align}
\mathcal{V}, \mathcal{Y} \subseteq \mathbb{R}^n
\end{align}
are subsets of $\mathbb{R}^n$.

Equivariant networks attempt to learn a function that is invariant to the action of some group $G$. That is, we seek to learn a function $\psi_\theta$ such that

\begin{align}
\psi_\theta(g \cdot x) &= g \cdot \psi_\theta(x)
\end{align}
Given a set of group representations \begin{align}
T_g: \mathcal{V} \to \mathcal{V} 
\end{align} for $g \in G$, then a function $\psi$ is equivariant to $G$ if there exists a transformation 
\begin{align}
 S_g: \mathcal{Y} \to \mathcal{Y}   
\end{align} such that
\begin{align}
S_g(\psi_\theta(v)) &= \psi_\theta(T_g(v))
\end{align}
In general given $(T_g, S_g)$ we want to learn $\psi_\theta$ from data. We say that such a $\psi_\theta$ is a $G$-equivariant neural network.

You will prove in the exercises that a composition of equivariant maps is equivariant, so if we can construct a suitable set of equivariant layers $e_i$, we can build rich families of equivariant architectures through composition
\begin{align}
e &= e_1 \circ e_2 \circ \dotsc \circ e_k
\end{align}

\section{Linear Layers through Equivariant Convolutions}

We will address the construction of linear and nonlinear equivariant layers separately since they require different transformations.

The linear portion of traditional convolutional layer we have defined previously is naturally equivariant to translations. (This true up to artifacts of padding and edge effects; to make it mathematically true, consider an infinitely large grid $\mathbb{Z}^2$ for the image in question).

We can define a general notion of a group convolution for group $G$ by the equation
\begin{align}
(\psi \ast \phi)(h) &= \sum_g \psi(g) \phi(g^{-1} h)
\end{align}
Here is a conversion of this basic definition into Physika
\begin{lstlisting}[language=python]
class DiscreteGroupConvolution(G: Group, w: $\mathbb{R}[|G|]$):
  def $\lambda$(X: $\mathbb{R}[|G|]$) -> $\mathbb{R}[|G|]$:
    X$'$ = zeros(|G|)
    for g h:
      X$'$[h] += w[g]*X[$g^{-1}$*h]
    return X$'$
\end{lstlisting}
If the group is continuous, the definition is given by
\begin{align}
(\psi \ast \phi)(h) &= \int_g \psi(g) \phi(g^{-1} h) d\mu
\end{align}
The integral is defined with respect to the "Haar measure", which provides a mathematically sound way of integrating across broad families of continuous groups. Constructing a general layer for a continuous group is trickier than for a discrete group since the infinite summation cannot be done directly and has to to be performed by some finite approximation (such as summing up terms in a fourier decomposition).

It turns out that such group convolutions provide a natural framework to construct $G$-equivariant layers. In fact, under broad conditions, a linear function is $G$-equivariant if and only if it is defined by a group convolution.

For finite groups, the group convolution equation can be directly implemented. For continuous groups, matters become more complicated.

\section{Representation Theory}

Group representations prove useful for modeling $G$-equivariant functions when $G$ is a continuous group.

Suppose that group $G$ has irreducible representations
\begin{align}
\rho_0,\dotsc, \rho_k
\end{align}
Suppose that we had functions
\begin{align}
f_0: \rho_0(G) \to \mathbb{R}^n,\dotsc, f_k: \rho_k(G) \to \mathbb{R}^n
\end{align}
that acted respectively on the sets of matrices output by the representations $\rho_i$. Suppose that $\rho_i(G) \subseteq \mathbb{R}^{n \times n}$ Then we can construct a function $f$ by "summing" these constituent functions
\begin{align}
f &: \rho_0(G) \oplus \dotsc \oplus \rho_k(G) \to \mathbb{R}^{kn} \\
f &= f_0 \oplus \dotsc \oplus f_k
\end{align}

For the specific case of linear functions, $f_i$ can be represented by a vector of components $w_i$ (since each $\rho_i(g)$ is a $n \times n$ matrix for any $g \in G$). In this case, $f$ is just a vector of parameters

\begin{align}
&f_i \in \mathbb{R}^{n} \\
&f \in \mathbb{R}^{kn}
\end{align}

Irreducible representations provide a sort of "basis set" for the group, so it is in fact possible to decompose any function on the group $G$ into functions on the irreducible representations. Consequently, if you can construct meaningful group convolutions for each of the irreducible representations, you can construct a group convolution for the full group. In future chapters, we will explore how to construct explicit equivariant layers for more complex groups.

\section{Equivariant Nonlinearities}

A challenge in construcing equivariant architectures is that we want to construct nonlinear operations that are equivariant. The Clebsch-Gordon coefficients we saw previously prove very useful for providing a method of blending information across the irreducible representations of a group
\begin{align}
f_i' &= \sum_j \sum_k CG_{j,k,i} f_i f_j
\end{align}

If we view the $f_i$ as vectors, then we can view this Clebsch-Gordon transformation as an equivariant nonlinear transformation.

%We say that $f$ is $G$-equivariant if
%\begin{align}
%    f(\rho(g)x) = \rho'(g)f(x)
%\end{align}
%Here $\rho$, $\rho'$ are possibly different representations. %Let's look at an example of a fluid flow.
%\begin{align}
%    f: \mathbb{R}^{H \times W \times 2} \to \mathbb{R}^{H \times W \times 2}\\
%    (x, v) \to (x, 2v)
%\end{align}

%How do you constrain a neural network to be equivariant? First, a composition of equivariant functions is equivariant. The challenge is to make linear and nonlinear maps equivariant with respect to the symmetry group $G$ in question.

%\begin{align}
%    f: \mathbb{R}^3 \to \mathbb{R}^3 \\
%    x \mapsto Mx
%\end{align}

%$C_3$ cyclicly permutes $\mathbb{R}^3$. Then we claim
%\begin{align}
%    M = \begin{pmatrix}
%    a & b & c \\
%    ....
%    \end{pmatrix}
%\end{align}
%How can we make the activation equivariant?

%Interested in simulating fluid flow
%\begin{align}
%    \nabla \cdot v &= 0 \\
%    \frac{\partial v}{\partial t} &= - (v \cdot \nabla) v - \frac{1}{\rho_0} \nabla p + \mu \nabla^2 v + f
%\end{align}
%Has symmetries from galilean invariance. That is independence There is also a scaling low
%\begin{align}
%    T^{sc}_\lambda w(x, t) &= \lambda w(\lambda x, \lambda^2 t), \lambda \in \mathbb{R}_{\geq 0}
%\end{align}
%The scaling is handled by a truncated G-convollution.
%\begin{align}
%    v(p) = \sum_{q \in \mathbb{Z}^2} w(p + q) K(q) \\
%\end{align}
%For scale equivariant convolutions we have to translate and scale kernal across the input
%\begin{align}
%    v(p, s, \mu) = \sum_{\lambda \in \mathbb{R}_{\geq 0},t\in\mathbb{R},q\in\mathbb{Z}^2} \lambda q(\lambda p + q, \lambda^2t, \lambda \mu) K(q, s, t, \lambda)
%\end{align}
%Uniform motion is handled by "conjugated canonicalization."

%E2-steerable CNN
%\begin{align}
%    T^{Rot}_R w(x, t) &= Rw(R^{-1}x, t), R \in O(2)\\
%\end{align}
%This is a $G$-equivariant network.
%\begin{align}
%    \rho_1 : G \to Gl(\mathbb{R}^2) \\
%    g \mapsto \begin{bmatrix}
%    \cos(\theta) & \sin(\theta) \\
%    -\sin(\theta) & \cos(\theta) 
%    \end{bmatrix}
%\end{align}
%Tested on Rayleigh-Benard convection.

\section{Exercises}
\begin{enumerate}
    \item Prove that the composition of two $G$-equivariant functions is $G$-equivariant.
\end{enumerate}
\end{document}